% !TEX program = xelatex
%----------------------- Преамбула -----------------------
\documentclass[14pt]{templateReportBMSTU}

% ---------------------- Документ ----------------------
\begin{document}
	\begin{titlepage}
		\noindent\begin{minipage}{0.05\textwidth}
			\includegraphics[scale=0.3]{1.png}
		\end{minipage}
		\hfill
		\begin{minipage}{0.85\textwidth}\raggedleft
			\begin{center}
				\fontsize{11pt}{0.4\baselineskip}\selectfont 
				\textbf{Министерство науки и высшего образования Российской Федерации
				\\ Федеральное государственное бюджетное образовательное учреждение
				\\ высшего образования
				\\ <<Московский государственный технический университет
				\\ имени Н.Э. Баумана
				\\ (национальный исследовательский университет)>>
				\\ (МГТУ им. Н.Э. Баумана)}
			\end{center}
		\end{minipage}

		\begin{center}
			\fontsize{11pt}{0.1\baselineskip}\selectfont
			\noindent\makebox[\linewidth]{\rule{\textwidth}{4pt}} \makebox[\linewidth]{\rule{\textwidth}{1pt}}
		\end{center}

		\begin{flushleft}
			\fontsize{11pt}{1\baselineskip}\selectfont 
			
			ФАКУЛЬТЕТ \textbf{ИНФОРМАТИКА И СИСТЕМЫ УПРАВЛЕНИЯ}

			КАФЕДРА \textbf{КОМПЬЮТЕРНЫЕ СИСТЕМЫ И СЕТИ (ИУ6)}

			НАПРАВЛЕНИЕ ПОДГОТОВКИ \textbf{09.XX.XX \ Название наравления}
		\end{flushleft}

		\begin{center}
			\hfill
			\\ \
			\fontsize{22pt}{\baselineskip}\selectfont
			\textbf{РАСЧЕТНО-ПОЯСНИТЕЛЬНАЯ ЗАПИСКА}
			\fontsize{20pt}{\baselineskip}\selectfont
			\\ \textbf{\textit{К \ \ КУРСОВОЙ \ \  РАБОТЕ}}
			\\ \textbf{\textit{по дисциплине <<XXXXXXXXXXX>>}}
			\\ \textbf{\textit{НА ТЕМУ:}}
			\\ \uline{\hfill}
			\\ \uline{\hfill}
			\\ \uline{\hfill}
			\\ \uline{\hfill}
		\end{center}

		\begin{table}[h!]
			\fontsize{12pt}{0.7\baselineskip}\selectfont
			\centering
			\begin{signstabular}[0.7]{p{3.25cm} >  {\centering\arraybackslash}p{4cm} > {\centering\arraybackslash}p{4cm} > {\centering\arraybackslash}p{4cm}}
				Студент & \uline{\hspace*{4cm}} & \uline{\hspace*{4cm}} & \uline{\hfill \textbf{<И.О. Фамилия>} \hfill} 
				\\ & \scriptsize (Группа) & \scriptsize (Подпись, дата) & \scriptsize (И.О. Фамилия)
			\end{signstabular}

			\vspace{\baselineskip}

			\begin{signstabular}[0.7]{p{3.25cm} >  {\centering\arraybackslash}p{4cm} > {\centering\arraybackslash}p{4cm} > {\centering\arraybackslash}p{4cm}}
				Руководитель & \ & \uline{\hspace*{4cm}} & \uline{\hfill \textbf{<И.О. Фамилия>} \hfill} 
				\\ & \ & \scriptsize (Подпись, дата) & \scriptsize (И.О. Фамилия)
			\end{signstabular}
		\end{table}

		\vfill
		\begin{center}
			\normalsize \textit{<Год> г.}
		\end{center}
	\end{titlepage}

	\normalsize
	\setcounter{page}{4}

	\pagebreak

	\section*{РЕФЕРАТ}
		\begin{center}
			Расчетно-пояснительная записка \pageref{LastPage} с., \totalfigures\ рис., \totaltables\ табл., Х ист., Х прил.

			\textbf{<Ключевые слова>}
		\end{center}

		\pagebreak

	% Переопределяем название \toc и выводим сам \toc
	\renewcommand{\contentsname}{\normalsize\bfseries\centering СОДЕРЖАНИЕ}
	\small
	\tableofcontents
	\normalsize

		\pagebreak

	\section*{ПЕРЕЧЕНЬ СОКРАЩЕНИЙ И ОБОЗНАЧЕНИЙ}

		\pagebreak

	\section*{ВВЕДЕНИЕ}
		\addcontentsline{toc}{section}{ВВЕДЕНИЕ}

		\pagebreak

	\section{<Раздел>}

	ПРИМЕР НУМЕРОВАННОГО СПИСКА:
	\begin{enumerate} 
  		\item первый элемент первого уровня содержит список 
		\item  второй элемент первого уровня
	\end{enumerate}
	Пример ссылки на раздел: Подробная информация в разделе \ref{section1}
	\\ ПРИМЕР МАРКИРОВАННОГО СПИСКА:
	\begin{itemize}
		\item Each list item is ...
		\item Lists can be nested ...
			\begin{itemize}
				\item The maximum ...
				\item Switching ...
			\end{itemize}
		\item And so on.
	\end{itemize}
	
	Листинг \thenumberlisting \ - Код программы
	\begin{lstlisting}
		c = x^2
		print('test listing')
	\end{lstlisting}
	\pagebreak

	\section{<Раздел>}  \label{section1}

	\begin{tabular}{ | p{100pt} | r | c | }
		\hline
		ИСЗ & Дата запуска & Масса, кг  \\ \hline
		Спутник-1 & 4 октября 1957 & 83.6 \\
		Спутник-2 & 3 ноября 1957 & 508.3 \\
		Эксплорер-1 & 1 февраля 1958 & 21.5 \\
		\hline
	\end{tabular}

	\begin{figure}[h]
		\centering
		\includegraphics[width=0.5\linewidth]{1.png}
		\caption{Тестовый рисунок}
		\label{fig:mpr}
	\end{figure}
	
	Формула: $c^{2}=a^{2}+b^{2}$
	\begin{math}
		c^2 = a^2 + b_1^2
	\end{math}
		\pagebreak
	
	\section{123}
		\pagebreak

	\section*{ЗАКЛЮЧЕНИЕ}
		\addcontentsline{toc}{section}{ЗАКЛЮЧЕНИЕ}

		\pagebreak

	\section*{СПИСОК ИСПОЛЬЗОВАННЫХ ИСТОЧНИКОВ}
		\addcontentsline{toc}{section}{СПИСОК ИСПОЛЬЗОВАННЫХ ИСТОЧНИКОВ}

		\pagebreak
\end{document}